% !TEX TS-program = xelatex
% !TEX encoding = UTF-8 Unicode
% !Mode:: "TeX:UTF-8"

\documentclass{resume}
\usepackage{zh_CN-Adobefonts_external}
\usepackage{linespacing_fix}
\usepackage{cite}

\begin{document}
\pagenumbering{gobble}

\name{詹勇}

\contactInfo{(+86) 188-5660-3434}{3446914219@qq.com}{GitHub | CSDN @daetz-coder}{浙江-杭州}

\section{个人简介}
计算机技术硕士在读，聚焦 \textbf{AI Agent 系统设计、RAG 与 LLM 应用工程化}。独立完成 Agent 运行时评估平台、RAG 知识库系统（Milvus 向量检索 + BM25 混合搜索）等全栈项目，熟悉 LangGraph 工作流编排、混合检索增强与 LLM-as-Judge 评估方法。

\section{教育背景}
\datedsubsection{\textbf{浙江理工大学}, 计算机技术硕士（在读）}{2024.9 - 至今}
\ 一等优秀奖学金(1次)，二等优秀奖学金(1次)
\datedsubsection{\textbf{嘉兴大学}, 软件工程学士}{2020.9 - 2024.6}
\ 校级优秀毕业生，浙江省政府奖学金(1次)，学习优秀奖学金(6次)

\section{学术成果}
\begin{itemize}[parsep=0.1ex,topsep=0ex]
  \item \textbf{Yong Zhan}, Xiangyang Li, et al. ``W-MSTC: A Window Based Multi-Scale Temporal Cross Attention Method for Few-Shot Modulation Recognition,'' \textit{IEEE Communications Letters}, 2025.（\textbf{第一作者}）
  \item Congyuan Xu, \textbf{Yong Zhan}, et al. ``Multimodal Fusion Based Few-Shot Network Intrusion Detection System,'' \textit{Scientific Reports}, 2025.（\textbf{第二作者}）
\end{itemize}

\section{技术能力}
\begin{itemize}[parsep=0.1ex,topsep=0ex]
  \item \textbf{Agent / RAG}: LangGraph、StateGraph、Human-in-the-Loop、Checkpoint、Tool Calling、LLM-as-Judge、Milvus 向量检索、BM25、BGE-M3、RRF、RecursiveCharacterTextSplitter
  \item \textbf{LLM 应用}: DeepSeek / GLM-4 / Qwen-Plus 多模型共识、Prompt/Context Engineering、SSE 流式输出、Token 优化
  \item \textbf{全栈工程}: Python / FastAPI / Async SQLAlchemy / Pydantic / Vue3 / TypeScript / ElementPlus / ECharts / Milvus / Docker
\end{itemize}

\section{项目经历}
\datedsubsection{\textbf{Agent Runtime Evaluation Platform：AI Agent 运行时质量评估平台} \quad \textit{Python / FastAPI / LangGraph / Vue3}}{2025.3 - 至今}
\begin{itemize}[parsep=0.15ex,topsep=0ex]
  \item \textbf{六维评估体系（17 项子指标 + 幻觉检测）：}Planning（覆盖度/顺序/粒度/完整性）、Tactical（相关性/效率/正确性）、Tool Use（选择质量/参数准确性/结果利用）、Memory（保持力/相关性/一致性）、Replan（触发适当性/适应质量/学习能力）、Retrieval（相关性/证据准确性/覆盖度 + hallucination\_detected），加权聚合 + 四档标定（$\geq$80/$\geq$60/$\geq$40 优秀至较差），6 条不同质量轨迹 $\times$ 6 评估器实测：综合分 \textbf{93.1 $\rightarrow$ 92.0 $\rightarrow$ 81.0 $\rightarrow$ 54.4 $\rightarrow$ 27.8 $\rightarrow$ 20.0} 单调递减，Planning 最敏感（0.0–80.0），并行评估耗时 \textbf{$\sim$15s}（asyncio.gather），成本 \textbf{¥0.012}（DeepSeek，对比 GPT-4o 低 27x）
  \item \textbf{多模型共识评估：}支持 DeepSeek + GLM-4 + Qwen-Plus 三厂商独立评分，输出 mean $\pm$ std（一致性指标），std $>$ 15 告警高分歧；提供 Consensus API + EvaluationDetail 前端多模型对比柱状图
  \item \textbf{三种零侵入接入模式：}LangGraph Instrument（node 函数包装 + 自动 tool\_calls 提取）、LLM Proxy（BaseChatModel 透明代理）、LangChain Callback（on\_llm/tool 生命周期钩子），均一行代码；独立 SDK 包（sdk/），外部项目仅需 httpx + langchain-core
\end{itemize}

\datedsubsection{\textbf{Wiki-Agent：AI 驱动的智能知识库系统} \quad \textit{LangGraph / FastAPI / Vue3 / Milvus}}{2025.5 - 至今}
\begin{itemize}[parsep=0.15ex,topsep=0ex]
  \item \textbf{多格式文档解析 + RecursiveCharacterTextSplitter 分块：}支持 Markdown/PDF/Word/TXT 自动识别与解析，LangChain RecursiveCharacterTextSplitter 分层分隔符（段落→句子→标点）分块（500 字符/50 重叠），替代手写 chunker，更稳健的语义边界保持
  \item \textbf{三级混合检索（20 条查询评估）：}Milvus 语义搜索（BGE-M3 + cosine）Top-1=\textbf{85.0\%}/MRR=\textbf{0.91}，BM25 关键词（jieba + rank\_bm25 + 25 词停用词表）Top-1=80.0\%/MRR=0.85，Hybrid RRF 融合（k=60）Top-1=85.0\%/MRR=0.87
  \item \textbf{4-Node StateGraph Agent：}search $\rightarrow$ respond $\rightarrow$ decide $\rightarrow$ execute，条件路由 + LangGraph interrupt 实现 Human-in-the-Loop CRUD 确认，AsyncSqliteSaver Checkpoint 中断恢复；PydanticOutputParser 驱动从对话自动提取知识（create/update/delete）
  \item \textbf{四路实时同步 + 向量管理后台：}Markdown $\rightarrow$ RecursiveCharacterTextSplitter 分块 $\rightarrow$ BGE-M3 Embedding $\rightarrow$ Milvus 向量库 + BM25 倒排索引 + Git 版本记录，增/改/删保持四路一致性，支持 commit 级回滚 + 一键全量重建；Vue3 + Element Plus 向量管理后台，可视化管理 Milvus 分块与元数据
\end{itemize}

\end{document}
